\renewcommand{\thetable}{3.A.5}
{\LFXcuerpotabla\setlength{\LTpre}{12pt}\setlength{\LTpost}{16pt}%
\begin{longtable}{@{}F{0.091} F{0.172} F{0.381} F{0.144} F{0.121} F{0.091}@{}}
\caption{Requerimientos no funcionales derivados de las Bases Técnicas Transversales. Fuente: elaboración propia a partir de las Bases Técnicas Transversales.}\label{tab:3-A-5}\\[2pt]
\toprule \cab{ID} & \cab{Materia} & \cab{Umbral o criterio} & \cab{Verificación prevista} & \cab{Prioridad / etapa} & \cab{Origen} \\ \midrule
\endfirsthead
\multicolumn{6}{@{}l@{}}{\fontsize{9pt}{10.5pt}\selectfont\color{lafroxGris}Tabla \thetable\ --- continuación}\\[3pt]
\toprule \cab{ID} & \cab{Materia} & \cab{Umbral o criterio} & \cab{Verificación prevista} & \cab{Prioridad / etapa} & \cab{Origen} \\ \midrule
\endhead
\midrule
\multicolumn{6}{@{}l@{}}{\fontsize{9pt}{10.5pt}\selectfont\color{lafroxGris}continúa en la página siguiente}\\
\endfoot
\bottomrule
\endlastfoot
  RNF-13.01 & Disponibilidad del sitio on-premise durante corte de enlace (24h) & El componente on-premise debe operar de forma autónoma y degradada ante la pérdida total del enlace con la nube, durante un período mínimo de 24 horas continuas, continuando registrando las transacciones operacionales críticas de forma local, con integridad garantizada y sin pérdida de datos. & Prueba de desconexión controlada y reconciliación & Crítica / E1 y E2 & BTT, RT-03.10, RT-03.11, RT-03.12, RT-03.13 \\
  RNF-13.02 & Declaración de funciones no disponibles en modo desconectado & El PROPONENTE debe declarar qué funciones NO estarán disponibles en modo desconectado y qué procedimiento manual las suple. & Inspección documental del entregable & Crítica / E1 y E2 & BTT, RT-03.13 \\
  RNF-13.03 & Redundancia de equipos on-premise críticos & Los equipos on-premise críticos deben ser redundantes. & Prueba funcional con datos representativos en marcha blanca & Alta / E1 y E2 & BTT, RT-03.14 \\
  RNF-13.04 & Tolerancia del almacenamiento local a la falla de un disco & El almacenamiento local debe tolerar la falla de al menos un disco & Prueba funcional con datos representativos en marcha blanca & Alta / E1 y E2 & BTT, RT-03.14 \\
  RNF-13.05 & Declaración y justificación del nivel RAID & el PROPONENTE debe declarar el nivel RAID escogido y justificarlo frente a las alternativas. & Inspección documental del entregable & Alta / E1 y E2 & BTT, RT-03.14 \\
  RNF-13.06 & Endurecimiento de sistemas on-premise según CIS Benchmarks & Los sistemas on-premise deben endurecerse conforme a los CIS Benchmarks aplicables, con gestión centralizada de parches y ventana de aplicación acordada con el CLIENTE. & Prueba funcional con datos representativos en marcha blanca & Alta / E1 y E2 & BTT, RT-03.15 \\
  RNF-13.07 & Enlace redundante entre on-premise y nube & El enlace entre el sitio on-premise y la nube debe ser redundante, con caminos físicos y proveedores distintos, y conmutación automática con tiempo de conmutación declarado ($\le$ 5 min). & Ejercicio de recuperación con medición de RTO/RPO & Alta / E1 y E2 & BTT, RT-03.17 \\
  RNF-13.08 & Dimensionamiento de ancho de banda por sitio & El PROPONENTE debe dimensionar el ancho de banda requerido por sitio, en régimen normal y en peak, y justificarlo con el cálculo de volumen de transacciones y de datos. & Inspección documental del entregable & Alta / E1 y E2 & BTT, RT-03.20 \\
  RNF-13.09 & Estudio de cobertura de red inalámbrica en sitios operacionales (incluye cámaras de frío) & El PROPONENTE debe realizar un estudio de cobertura de red inalámbrica en los centros de distribución, que incluya expresamente el interior de las cámaras de refrigerado y de congelado, hoy sin señal, y proponer la solución técnica para resolverlo. & Inspección documental del entregable & Alta / E1 y E2 & BTT, RT-03.23 \\
  RNF-14.01 & Cifrado en tránsito con TLS 1.3 & Todo el tráfico en tránsito debe emplear TLS 1.3, con prohibición expresa de TLS 1.0 y 1.1, conjuntos de cifrado modernos, HSTS con precarga y gestión automatizada de certificados con rotación y alerta anticipada de vencimiento. & Prueba de desempeño con volumen de peak & Crítica / E1 y E2 & BTT, RT-11.08 \\
  RNF-14.02 & Cifrado en reposo de todos los datos & La totalidad de los datos en reposo debe estar cifrada, con claves gestionadas en un servicio de gestión de claves o en un módulo de seguridad de hardware, política de rotación declarada y separación de funciones en la custodia de claves. & Revisión de configuración y prueba de seguridad & Crítica / E1 y E2 & BTT, RT-11.09 \\
  RNF-14.03 & Segregación de red en la nube (subredes privadas) & En la nube, la red debe segmentarse por capas, con subredes privadas para aplicación y datos, y exposición pública restringida a la capa de borde. Ningún componente de datos debe ser alcanzable desde Internet. & Prueba funcional con datos representativos en marcha blanca & Alta / E1 y E2 & BTT, RT-03.04 \\
  RNF-14.04 & SIEM con casos de uso de detección específicos del negocio & Los eventos de seguridad deben registrarse de forma centralizada e inalterable, y correlacionarse en una plataforma SIEM con casos de uso de detección definidos específicamente para el proceso de negocio del caso (ej. intentos de acceso a trazabilidad de lote, modificaciones de inventario no autorizadas). & Revisión de configuración y prueba de seguridad & Alta / E1 y E2 & BTT, RT-11.14, RT-11.15 \\
  RNF-14.05 & Detección y respuesta en puntos finales (EDR) & Se debe implementar detección y respuesta en puntos finales (EDR) y en cargas de trabajo, tanto en nube como on-premise, con cobertura 24$\times$7. & Prueba de desempeño con volumen de peak & Alta / E1 y E2 & BTT, RT-11.16 \\
  RNF-14.06 & Pruebas de intrusión anuales por tercero independiente & Se deben ejecutar pruebas de intrusión (pentesting) por un tercero independiente del ADJUDICATARIO, anualmente y antes de cada paso a producción, con entrega íntegra del informe al CLIENTE y plan de remediación con plazos. LafroX ejecuta estas pruebas de forma semestral, por sobre el mínimo anual (LafroX, 2026b, sección «Modelo de Gobierno de Seguridad de la Información»). & Inspección documental del entregable & Alta / E1 y E2 & BTT, RT-11.20 \\
  RNF-14.07 & Análisis estático de código, SCA y escaneo de contenedores en CI/CD & El flujo de integración continua debe incorporar análisis estático de código (SAST), análisis de composición de software (SCA), análisis dinámico (DAST) y escaneo de imágenes de contenedor, con criterios de bloqueo automático del despliegue ante hallazgos críticos. & Ejercicio de recuperación con medición de RTO/RPO & Alta / E1 y E2 & BTT, RT-11.22 \\
  RNF-14.08 & Inventario de componentes de software (SBOM) en formato CycloneDX/SPDX & Cada versión liberada debe acompañarse de su inventario de componentes de software en formato CycloneDX o SPDX, entregado al CLIENTE. & Prueba funcional con datos representativos en marcha blanca & Alta / E1 y E2 & BTT, RT-11.23 \\
  RNF-14.09 & Verificación de procedencia de artefactos según SLSA nivel 3 & Los artefactos deben firmarse y su procedencia debe verificarse conforme a SLSA nivel 3 o superior. & Prueba funcional con datos representativos en marcha blanca & Alta / E1 y E2 & BTT, RT-11.24 \\
  RNF-14.10 & Prohibición de datos productivos en ambientes no productivos & Queda prohibido el uso de datos productivos reales en ambientes no productivos sin anonimización o seudonimización verificable. & Prueba funcional con datos representativos en marcha blanca & Alta / E1 y E2 & BTT, RT-11.25 \\
  RNF-15.01 & Política de sesión: duración, inactividad, renovación, revocación & El sistema debe implementar política de sesión con duración máxima, caducidad por inactividad, renovación de la credencial de sesión tras la autenticación, revocación inmediata y control de sesiones concurrentes. Los valores de duración, inactividad y concurrencia se declaran en la arquitectura de la oferta y se verifican antes del cierre de la marcha blanca E1. & Revisión de configuración y prueba de seguridad & Crítica / E1 y E2 & BTT, RT-12.07 \\
  RNF-15.02 & Credenciales de sesión firmadas y de vida breve & Las credenciales de sesión deben ser firmadas y de vida breve, con credencial de refresco rotatoria. Queda prohibido transportar identificadores de sesión en la ruta de la dirección web (URL). La vida de cada credencial se declara en la arquitectura de la oferta y se verifica antes del cierre de la marcha blanca E1. & Revisión de configuración y prueba de seguridad & Crítica / E1 y E2 & BTT, RT-12.08 \\
  RNF-15.03 & Registro completo del ciclo de vida de la identidad (auditoría) & El sistema debe registrar auditoría completa del ciclo de vida de la identidad: creación, modificación, elevación, bloqueo y baja de cuentas, con no repudio y retención declarada. & Inspección de política de retención y restauración de muestra & Alta / E1 y E2 & BTT, RT-12.09 \\
  RNF-16.01 & Indicadores de nivel de servicio medidos sobre experiencia real del usuario & Los indicadores de nivel de servicio deben medirse sobre la experiencia real de la persona usuaria (percentil 95) y no sobre pruebas sintéticas, sin perjuicio de que estas se empleen como complemento. & Prueba de desempeño con volumen de peak & Alta / E1 y E2 & BTT, RT-14.03 \\
  RNF-16.02 & Libro de operación y guía de resolución documentados & El PROPONENTE debe documentar un libro de operación y una guía de resolución para cada escenario de falla previsible, con automatización progresiva de las tareas repetitivas. & Inspección documental del entregable & Alta / E1 y E2 & BTT, RT-14.05 \\
  RNF-16.03 & Registros sin datos personales sensibles ni credenciales & Los registros (logs) de la solución no deben contener datos personales sensibles ni credenciales, y su acceso debe estar controlado y auditado. & Revisión de configuración y prueba de seguridad & Crítica / E1 y E2 & BTT, RT-14.07 \\
  RNF-17.01 & Retención de auditoría mínima de 5 años & El período de retención de la auditoría debe ser el que fije el caso y, en su defecto, no inferior a cinco años. & Inspección de política de retención y restauración de muestra & Alta / E1 y E2 & BTT, RT-16.10 \\
  RNF-17.02 & Exportaciones de gran volumen procesadas de forma asíncrona & Las exportaciones de gran volumen deben procesarse de forma asíncrona, con notificación al completarse y sin bloquear la sesión del usuario. & Revisión de configuración y prueba de seguridad & Alta / E1 y E2 & BTT, RT-16.29 \\
  RNF-17.03 & Portal público resistente a picos de tráfico & El portal público debe resistir picos de tráfico sin degradar los servicios transaccionales internos, mediante aislamiento de recursos y caché. & Prueba funcional con datos representativos en marcha blanca & Alta / E1 y E2 & BTT, RT-16.34 \\
  RNF-19.01 & Cálculo de capacidad con supuestos de usuarios concurrentes y TPS & El PROPONENTE debe presentar el cálculo de capacidad que sustenta su dimensionamiento, con los supuestos de usuarios concurrentes, transacciones por segundo, volumen de datos y crecimiento anual, tomados de la volumetría del caso. & Inspección documental del entregable & Alta / E1 y E2 & BTT, RT-09.01 \\
  RNF-19.02 & Escalamiento horizontal automático de capas de aplicación e integración & Las capas de aplicación e integración deben escalar horizontalmente de forma automática, con umbrales, límites superiores y costo asociado declarados en la oferta (arquitectura y Oferta Económica). & Prueba de carga que active el escalamiento en marcha blanca & Alta / E1 y E2 & BTT, RT-02.10, RT-09.04 \\
  RNF-19.03 & Identificación del cuello de botella al crecer la carga & El PROPONENTE debe identificar el componente que primero se convertirá en cuello de botella al crecer la carga y explicar cómo lo detectará y cómo lo resolverá. & Inspección documental del entregable & Alta / E1 y E2 & BTT, RT-09.05 \\
  RNF-19.04 & Pruebas de carga en Preproducción con 1.5$\times$ peak y pruebas de estrés & Se deben ejecutar pruebas de carga sobre Preproducción con un volumen equivalente a 1,5 veces el peak declarado, y pruebas de estrés hasta identificar el punto de quiebre de la solución. & Prueba de desempeño con volumen de peak & Alta / E1 y E2 & BTT, RT-09.06, RT-09.07 \\
  RNF-20.01 & Disponibilidad mensual mínima del 99.9\% para servicios críticos & La solución debe alcanzar una disponibilidad mensual mínima de 99.9\% para los servicios clasificados como críticos, medida sobre la transacción de negocio de extremo a extremo. & Ejercicio de recuperación con medición de RTO/RPO & Crítica / E1 y E2 & BTT, RT-10.01 \\
  RNF-20.02 & Clasificación de servicios por criticidad (crítico, alto, medio, bajo) & El PROPONENTE debe clasificar cada servicio de la solución en crítico, alto, medio o bajo, justificando la clasificación con el impacto operacional de su indisponibilidad, y aplicar a cada uno el nivel de servicio correspondiente. & Inspección documental del entregable & Crítica / E1 y E2 & BTT, RT-10.02 \\
  RNF-20.03 & Plan de continuidad del negocio conforme a ISO 22301 & El PROPONENTE debe elaborar un plan de continuidad del negocio conforme a ISO 22301, con análisis de impacto en el negocio, escenarios de contingencia, procedimientos manuales de respaldo y criterios de activación. & Inspección documental del entregable & Crítica / E1 y E2 & BTT, RT-10.03 \\
  RNF-20.04 & Pruebas de resiliencia con inyección de fallas & Se deben ejecutar pruebas de resiliencia mediante inyección controlada de fallas (caída de instancia, de zona, de dependencia externa, latencia elevada, saturación de disco) antes de cada paso a producción y al menos una vez por semestre durante la Operación. & Prueba de desempeño con volumen de peak & Alta / E1 y E2 & BTT, RT-10.07 \\
  RNF-20.05 & Documentación de comportamiento ante falla de dependencias externas & El PROPONENTE debe documentar, por cada dependencia externa, el comportamiento de la solución cuando esa dependencia no responde, responde con error o responde con lentitud. & Inspección documental del entregable & Alta / E1 y E2 & BTT, RT-10.08 \\
  RNF-20.06 & Site secundario para DR: RTO $\le$ 4h, RPO $\le$ 15 min & El PROPONENTE debe habilitar un sitio secundario para recuperación ante desastres, con RTO $\le$ 4 horas y RPO $\le$ 15 minutos para los servicios críticos, y probar la conmutación al menos 2 veces al año. & Inspección documental del entregable & Crítica / E1 y E2 & BTT, RT-07.04, RT-07.07 \\
  RNF-20.07 & Política de respaldo 3-2-1-1-0 (3 copias, 2 medios, 1 fuera de sitio, 1 inmutable, 0 errores) & La política de respaldo debe seguir el esquema 3-2-1-1-0: tres copias, en dos medios distintos, una fuera de sitio, una inmutable o fuera de línea y cero errores de verificación de restauración. & Ejercicio de recuperación con medición de RTO/RPO & Crítica / E1 y E2 & BTT, RT-07.09, RT-07.12 \\
  RNF-21.01 & Centro de operaciones de red (NOC) 24$\times$7$\times$365 & El ADJUDICATARIO debe disponer de un centro de operaciones de red (NOC) con cobertura 24$\times$7$\times$365, propio o subcontratado, y declarar su ubicación, dotación por turno y procedimientos. & Inspección documental del entregable & Crítica / E1 y E2 & BTT, RT-21.01 \\
  RNF-21.02 & Gerente de servicio dedicado para el CLIENTE & El ADJUDICATARIO debe designar un gerente de servicio dedicado como contraparte permanente del CLIENTE durante la fase de Operación. & Prueba funcional con datos representativos en marcha blanca & Alta / E1 y E2 & BTT, RT-21.02 \\
  RNF-21.03 & Centro de atención telefónica: 80\% en 20 segundos, 70\% resolución primer contacto & El centro de atención debe cumplir un tiempo de atención al usuario final de 80\% antes de 20 segundos y una resolución al primer contacto de al menos 70\%. La tasa de abandono no debe superar el 5\%. & Prueba de desempeño con volumen de peak & Alta / E1 y E2 & BTT, RT-21.06 \\
  RNF-21.04 & Horario de atención ampliado a 24$\times$7 para incidentes críticos y ventana de despacho (05:30-07:00) & El horario mínimo de atención será de 08:00 a 20:00 en días hábiles, ampliado a 24$\times$7 para los incidentes de severidad crítica y para la ventana operacional crítica (despacho de 05:30 a 07:00). & Prueba funcional con datos representativos en marcha blanca & Alta / E1 y E2 & BTT, RT-21.07 \\
  RNF-21.05 & Dimensionamiento de la dotación del centro de atención con teoría de colas o Erlang C & El PROPONENTE debe dimensionar la dotación del centro de atención con fundamento cuantitativo, empleando teoría de colas o el modelo Erlang C, a partir del volumen de contactos proyectado del caso. & Inspección documental del entregable & Alta / E1 y E2 & BTT, RT-21.14 \\
  RNF-21.06 & Canal único de registro de incidentes con ticket y seguimiento & Debe existir un canal único de registro de incidentes y solicitudes, con número de ticket, clasificación por severidad y seguimiento del ciclo de vida completo hasta el cierre conforme. & Prueba funcional con datos representativos en marcha blanca & Alta / E1 y E2 & BTT, RT-21.15 \\
  RNF-21.07 & Traslado a sitios alejados (Curicó, Chillán, Los Ángeles) para resolución de incidentes & Los especialistas de niveles 2 y 3 deben trasladarse a los sitios alejados (Curicó, Chillán, Los Ángeles) cuando la resolución lo requiera, por el medio más rápido disponible, sin alterar la atención normal del resto de los sitios. El costo del traslado está incluido en la oferta. & Ejercicio de recuperación con medición de RTO/RPO & Alta / E1 y E2 & BTT, RT-21.16 \\
  RNF-21.08 & Bolsa anual de horas de mantención evolutiva & El PROPONENTE debe declarar el tamaño de la bolsa anual de horas de mantención evolutiva, su composición por perfil y el procedimiento de solicitud, estimación, aprobación y liquidación. & Inspección documental del entregable & Alta / E1 y E2 & BTT, RT-21.19 \\
  RNF-21.09 & Planificación anual de actualizaciones de versiones de componentes & Las actualizaciones de versión de los componentes de base (sistemas operativos, bases de datos, frameworks) deben planificarse anualmente, con ventana acordada y plan de reversión. & Inspección documental del entregable & Alta / E1 y E2 & BTT, RT-21.22 \\
  RNF-22.01 & Plan de capacitación estructurado por perfil (usuarios, administradores, técnicos) & El PROPONENTE debe presentar un plan de capacitación estructurado por perfil: personas usuarias finales, usuarias avanzadas, administradoras, equipo técnico y soporte de niveles 1 y 2. & Inspección documental del entregable & Alta / E1 y E2 & BTT, RT-22.01 \\
  RNF-22.02 & Material de capacitación en español, editable y de propiedad del CLIENTE & Todo el material de capacitación debe entregarse en español, en formato editable (Word, PowerPoint, etc.) y de propiedad del CLIENTE: manuales por perfil, guías rápidas, preguntas frecuentes, videos tutoriales y base de conocimiento consultable. & Ejercicio de recuperación con medición de RTO/RPO & Alta / E1 y E2 & BTT, RT-22.03 \\
  RNF-22.03 & Capacitación sin afectar la operación: programación por turnos y horarios acordados & La capacitación no podrá afectar la operación del CLIENTE: se programará por turnos y en horarios acordados, considerando la estacionalidad (peaks de septiembre y diciembre) y la ventana operacional (despacho de mañana, preparación nocturna). & Prueba funcional con datos representativos en marcha blanca & Alta / E1 y E2 & BTT, RT-22.04 \\
  RNF-22.04 & Evaluación y certificación de usuarios administradores y equipo técnico & Las personas usuarias administradoras y el equipo técnico del CLIENTE deben ser evaluados y certificados como condición para el cierre de cada marcha blanca. & Prueba funcional con datos representativos en marcha blanca & Alta / E1 y E2 & BTT, RT-22.05 \\
  RNF-22.05 & Capacitación de personal nuevo del CLIENTE sin costo adicional durante la Operación & Durante la Operación, el ADJUDICATARIO debe ejecutar al menos dos jornadas anuales de actualización y capacitar al personal nuevo del CLIENTE sin costo adicional. & Prueba funcional con datos representativos en marcha blanca & Alta / E1 y E2 & BTT, RT-22.06 \\
  RNF-22.06 & Acompañamiento y mentoría al equipo técnico del CLIENTE durante 6 meses post producción & El ADJUDICATARIO debe proveer acompañamiento y mentoría al equipo técnico del CLIENTE durante al menos seis meses posteriores al paso a producción de la Etapa 2. & Prueba funcional con datos representativos en marcha blanca & Alta / E1 y E2 & BTT, RT-22.07 \\
  RNF-23.01 & Especificación de hardware de terreno con marca, modelo, cantidad y características & El PROPONENTE debe especificar cada dispositivo de terreno (lectores de código, impresoras térmicas, sensores de temperatura, terminales móviles) con marca, modelo de referencia, cantidad, características mínimas, accesorios, consumibles y costo unitario estimado. & Inspección documental del entregable & Alta / E1 y E2 & BTT, RT-08.10, RT-08.11 \\
  RNF-23.02 & Dispositivos con grado de protección IP y resistencia a caídas coherentes con el entorno & Los dispositivos de terreno deben declarar su grado de protección contra polvo y agua (IP) y su resistencia a caídas, coherentes con el entorno de operación (intemperie, humedad, polvo, vibración, temperatura). El grado IP y la altura de caída de cada modelo se declaran en la especificación de hardware de la oferta (RNF-23.01) y se verifican con la ficha del fabricante. & Inspección documental del entregable & Alta / E1 y E2 & BTT, RT-08.12 \\
  RNF-23.03 & Ciclo de vida esperado de cada dispositivo y plan de reposición durante 56 meses & El PROPONENTE debe indicar el ciclo de vida esperado de cada dispositivo de terreno, la disponibilidad de repuestos y el plan de reposición durante los 56 meses del Contrato. & Inspección documental del entregable & Alta / E1 y E2 & BTT, RT-08.13 \\
  RNF-23.04 & Estaciones de trabajo para operación y administración con monitores duales & El PROPONENTE debe especificar las estaciones de trabajo requeridas para la operación y administración de la plataforma, en la cantidad que determine el dimensionamiento, con monitores duales y gestión centralizada con cifrado de disco. Los puestos cumplirán las condiciones ergonómicas exigidas en RT-08.08, que remite a la NCh 2527 (\textit{Bases Técnicas Transversales}, 2026, cap. 8, RT-08.08, p. 19). & Inspección documental del entregable & Alta / E1 y E2 & BTT, RT-08.07, RT-08.08, RT-08.09 \\
\end{longtable}}
